%! TEX root = ../main.tex
\documentclass[../main.tex]{subfiles}

\begin{document}
\section{Power Electronics System}
\label{sec:power}

To enable physically untethered operation, we used an iterative process to engineer a fully power autonomous robotic system with 3 main components: a flying robot, a custom lithium-ion battery, and a power-electronics circuit. In this process, we co-designed each system to maximize the total power-density of the robot. This thesis focuses on the power-electronics circuit to implement a sub-gram high-voltage power electronics circuit using off-the-shelf components and commercial flex PCB manufacturing process. In the process, we investigated different converter implementations and driving waveforms to minimize battery weight and maximize mechanical output. The success of this project will unshackle the robot from stationary power amplifiers, bringing soft-actuated flying robots closer to practical applications.

\subsection{Battery Design}

\subsection{Waveform Optimization}

\subsection{Converter Optimization}

\subsection{Driver Circuits}

\subsection{Lightweight PCB}

\subsection{Premilinary Results}

\end{document}
